% !TEX program = pdflatex
\documentclass[11pt]{article}
\usepackage[margin=1in]{geometry}
\usepackage{amsmath,amssymb}
\usepackage{booktabs}
\usepackage{hyperref}
\usepackage{natbib}
\usepackage{listings}
\usepackage[T1]{fontenc}
\usepackage[utf8]{inputenc}
\usepackage{enumitem}

\hypersetup{colorlinks=true,linkcolor=black,citecolor=black,urlcolor=black}

\lstdefinestyle{code}{
  basicstyle=\ttfamily\small,
  frame=single,
  breaklines=true,
  numbers=left,
  numberstyle=\tiny,
  tabsize=2,
  showstringspaces=false
}
\lstset{style=code}

\title{Healthspan T2D Early--Risk Prediction: A Minimal, Clinically Grounded Protocol}
\author{Ryan O. van Gelder \& Ruth Russel \\ \\ Citizen Gardens}
\date{October 23, 2025}

\begin{document}
\maketitle

\begin{abstract}
We specify a pragmatic, publish--grade protocol to predict incident type 2 diabetes (T2D) at multi--horizon H=\{1,3,5\} years from routine longitudinal electronic health record (EHR) data. The design emphasises calibration, decision--analytic utility, fairness to baselines, and reproducibility. We include acceptance gates that prevent deployment without clinically meaningful benefit.
\end{abstract}

\section{Objective}
Predict incident T2D within H\,$\in$\,\{1,3,5\} years from routine labs and vitals. \textbf{Primary metric:} AUROC at H=3y. \textbf{Secondary:} AUROC at H=1y and 5y, AUPRC, expected calibration error (ECE), Brier score \citep{brier1950verification}, and decision--curve net benefit \citep{vickers2006decision}. \textbf{Success:} $\Delta$AUROC$\ge0.02$ versus XGBoost \citep{chen2016xgboost} and ECE$\le0.03$ on the held--out test set.

\section{Endpoint}
\textbf{Outcome (binary, each horizon scored independently):} first occurrence of any: (i) ICD--10 E11.* recorded twice on different days; (ii) HbA1c$\ge$6.5\% confirmed on a separate date; (iii) new prescription for a non--insulin antihyperglycemic agent. Index date $t_0$ is the last encounter in the observation window. Exclude prior evidence of T2D before $t_0$.

\section{Cohort}
Adults 18--90 with $\ge$3 encounters and $\ge$12 months of history before $t_0$. Target $N\ge10{,}000$ and event rate $\ge10\%$ or class--balanced sampling. Person--level 70/15/15 train/validation/test with temporal separation. If a second site exists, external validation is required; otherwise construct a strict temporal external test using the most recent 12 months.

\section{Data and Features}
Structured EHR only: demographics; vitals; labs (HbA1c, fasting glucose, lipid panel, ALT/AST, creatinine/eGFR, CRP, CBC); meds (antihypertensives, statins, steroids, atypical antipsychotics); utilisation features. Use the last 24 months [$-T,0$); include binary observed/missing flags; no forward--fill beyond 90 days.

\section{Preprocessing}
Monthly binning; per--site standardisation using training statistics; clip at 0.5th/99.5th percentiles; encode meds as monthly exposure counts.

\section{Baselines}
\begin{enumerate}[label=\arabic*)]
  \item Logistic regression (L2) on engineered features: last, mean, and slope over 6/12/24 months.
  \item XGBoost \citep{chen2016xgboost} on the same engineered features.
  \item Feature--hashing logistic regression \citep{weinberger2009feature} on sparse event codes (2$^i$ buckets, $i\in\{12,13,14\}$).
  \item Clinician rule (positive control): risk=1 if last HbA1c$>$6.0\% or fasting glucose$\ge$110 mg/dL; else 0.
\end{enumerate}

\section{Proposed Model}
Sequence classifier on monthly slices. For patient $i$ with sequence $X_i=\{x_{i,1},\ldots,x_{i,T}\}$ and $x_{i,t}\in\mathbb{R}^d$, embed $z_{i,t}=Wx_{i,t}+b+e_t$ with learnable positional vectors $e_t$. Backbone: 2--4 layer Transformer encoder \citep{vaswani2017attention} with masking for missing months. Head: $h_i=\mathrm{Pool}(\{z'_{i,t}\})$; risk $\hat p_i=\sigma(w^\top h_i + c)$. Regularisation: dropout 0.1--0.3; weight decay $10^{-5}$.

\paragraph{Fairness to baselines.} Append summary features (last, mean, slope over 6/12/24 months) to the sequence representation before the head.

\paragraph{Multi--horizon training.} Single shared encoder with horizon--specific heads. Add a horizon token $e_H$ to each timestep embedding. One sigmoid output per $H\in\{1,3,5\}$. Overall loss $\mathcal{L}=\sum_H \alpha_H\,\mathrm{BCE}_H$ with $\alpha_H\propto 1/\mathrm{prevalence}_H$ and $\sum_H \alpha_H=1$. Early stopping monitors H=3y AUROC; all horizons are reported.

\paragraph{Ablations.} Replace Transformer with GRU; remove time embeddings; depth/hidden--size sweep; prime--hashing variant mapping high--cardinality IDs via $b=(a\cdot\mathrm{ID}+c)\bmod P$ compared to standard feature hashing \citep{weinberger2009feature}.

\section{Training Protocol}
Weighted binary cross--entropy with class weights from training prevalence; focal--loss sensitivity \citep{lin2017focal}. Optimiser AdamW \citep{loshchilov2019adamw} with lr in \{1e--4, 3e--4\}. Early stopping on validation AUROC (H=3y) with patience 10 epochs. Batch by patient, pad months, mask missing. Calibrate predictions using Platt scaling \citep{platt1999probabilistic} or isotonic regression \citep{zadrozny2002transforming} trained on validation only. Validation/test remain at natural prevalence; report PR curves and prevalence--adjusted PPV/NPV. Compute budget: $\le$1\,$\times$\,A100 40GB or $\le$2\,$\times$\,RTX 3090 for $<$24 hours total across all runs including ablations.

\section{Evaluation}
Report on the test set once. AUROC and AUPRC with 1{,}000 bootstrap CIs \citep{efron1979bootstrap}; ECE (10 bins); Brier score \citep{brier1950verification}; calibration plot and reliability table; decision--curve analysis \citep{vickers2006decision} at thresholds spanning 1--20\%; subgroup performance by sex, age bands, and race/ethnicity. Temporal robustness: evaluate per calendar quarter, train on [$-T,0$) and test on $(0,T']$ with $T'$ the most recent 12 months. 

\paragraph{Risk persistence rule.} ``High--risk'' only if risk$\ge\tau$ at $\ge2$ encounters within $\le6$ months; report change in PPV, sensitivity, and alert volume vs single--encounter rule. Optional stability penalty ablation adds $\lambda\,\mathrm{mean}_t\,|p_{t+1}-p_t|$ with $\lambda\in\{0,0.01,0.05\}$ and reports calibration.

\paragraph{Clinical utility.} Number Needed to Screen (NNS) at $\tau$; lead--time gain versus standard care; PPV in predefined high--risk subgroups; standard and standardised net benefit \citep{vickers2006decision}.

\paragraph{Interpretability.} SHAP for XGBoost \citep{lundberg2017unified}; Integrated Gradients for the sequence model \citep{sundararajan2017axiomatic}; rank top--10 contributing features per subgroup; clinical panel face--validity ratings and adjudication notes.

\paragraph{Error analysis.} Manual chart review of $\ge50$ high--confidence false positives and $\ge50$ false negatives to categorise causes (label noise, data gaps, physiology). Publish taxonomy and remediation actions.

\paragraph{Statistical tests.} DeLong test for AUROC differences \citep{delong1988comparing}; McNemar's test at the chosen threshold \citep{mcnemar1947note}.

\section{Reproducibility and Governance}
Pre--registration; containerised training with deterministic seeds; full code release with configs, feature dictionary, and model card; versioned data processing lineage. PHI never leaves site; multi--site training uses federated aggregation \citep{kairouz2021advances}.

\section{Risks and Mitigations}
\textbf{Recalibration/rollback triggers:} ECE$>$0.05 for two consecutive quarters; or AUROC drop $\ge0.03$ for two consecutive quarters; or Spiegelhalter's calibration $z$ test $p<0.01$ \citep{spiegelhalter1986probabilistic}; or PSI$>$0.2 on any top--10 feature. Actions: recalibrate on last 12 months; if unresolved, retrain; if unresolved, rollback.

Other risks: label noise (confirmatory rules; Rx--only removal sensitivity), selection bias (flow diagram; included vs excluded comparison), overfitting (baselines, ablations, early stopping), data quality (plausibility checks; unit harmonisation).

\section{Acceptance Gate}
Proceed only if: (i) H=3y: $\Delta$AUROC$\ge0.02$ vs XGBoost \emph{and} ECE$\le0.03$ on test \emph{and} decision--curve shows positive net benefit at a clinically chosen $\tau$; (ii) H=1y and H=5y show no clinically material degradation when evaluated per--horizon, else they are reported as negative results; (iii) under the persistence rule at $\tau^*$, PPV improves with $\le25\%$ relative loss in sensitivity vs single--encounter rule. Otherwise stop and publish the negative result.

\section{Repository Skeleton and Code Snippets}
\subsection{Layout}
\begin{lstlisting}
healthspan-t2d/
  data_contract/columns.md
  sql/cohort.sql
  configs/default.yaml
  src/models/transformer.py
  src/train.py
  src/eval.py
  src/utils/metrics.py
  notebooks/calibration.ipynb
  reports/model_card.md
\end{lstlisting}

\subsection{configs/default.yaml}
\begin{lstlisting}[language=bash]
seed: 13
horizons: [1,3,5]
window_months: 24
batch_size: 256
lr: 0.0003
dropout: 0.2
weight_decay: 1e-5
transformer:
  layers: 3
  d_model: 128
  n_heads: 4
  ff_mult: 4
loss:
  type: bce
  alpha_by_prevalence: true
early_stopping:
  monitor: auroc_h3
  patience: 10
calibration: isotonic
\end{lstlisting}

\subsection{src/models/transformer.py}
\begin{lstlisting}[language=Python]
import torch
import torch.nn as nn

class MHTransformer(nn.Module):
    """Multi-horizon Transformer with shared encoder and horizon-specific heads."""
    def __init__(self, d_in, d_model=128, layers=3, n_heads=4, horizons=(1,3,5), max_months=24):
        super().__init__()
        self.horizons = [str(h) for h in horizons]
        self.proj = nn.Linear(d_in, d_model)
        self.pos = nn.Embedding(max_months + 1, d_model)  # months + pad
        self.horiz = nn.Embedding(8, d_model)             # horizon token index
        enc_layer = nn.TransformerEncoderLayer(d_model, n_heads, d_model*4, batch_first=True)
        self.enc = nn.TransformerEncoder(enc_layer, layers)
        self.heads = nn.ModuleDict({h: nn.Linear(d_model, 1) for h in self.horizons})

    def forward(self, x, t_idx, h_idx, mask=None, return_logits=True):
        # x: [B, T, d_in], t_idx: [B, T], h_idx: [B, T], mask: [B, T] (True for PAD)
        z = self.proj(x) + self.pos(t_idx) + self.horiz(h_idx)
        z = self.enc(z, src_key_padding_mask=mask)
        # mean pool over valid timesteps
        denom = (~mask).sum(1).clamp_min(1).unsqueeze(-1)
        h = (z.masked_fill(mask.unsqueeze(-1), 0).sum(1)) / denom
        out = {k: self.heads[k](h).squeeze(-1) for k in self.heads}  # logits
        if return_logits:
            return out
        return {k: torch.sigmoid(v) for k, v in out.items()}
\end{lstlisting}

\subsection{src/utils/metrics.py}
\begin{lstlisting}[language=Python]
from sklearn.metrics import roc_auc_score, average_precision_score, brier_score_loss

def auroc(y, p):
    return roc_auc_score(y, p)

def auprc(y, p):
    return average_precision_score(y, p)

def brier(y, p):
    return brier_score_loss(y, p)
\end{lstlisting}

\subsection{src/train.py (sketch)}
\begin{lstlisting}[language=Python]
import torch, torch.nn as nn
from torch.utils.data import DataLoader

# dataset yields (x, t_idx, h_idx, mask, labels_dict) where labels_dict['1'],['3'],['5']

bce = nn.BCEWithLogitsLoss(reduction='none')
alpha = {1: 0.5, 3: 0.3, 5: 0.2}  # set ∝ 1/prevalence then normalize

for epoch in range(max_epochs):
    model.train()
    for x, t_idx, h_idx, mask, y in DataLoader(train_ds, batch_size=cfg.batch_size, shuffle=True):
        out = model(x, t_idx, h_idx, mask, return_logits=True)
        loss = 0.0
        for H in (1,3,5):
            logits = out[str(H)]
            target = y[str(H)].float()
            w = class_weights[str(H)]  # tensor scalar per-H
            l = bce(logits, target)
            l = (w*target + (1-target)) * l  # pos_weight-like effect
            loss = loss + alpha[H] * l.mean()
        opt.zero_grad(); loss.backward(); opt.step()
\end{lstlisting}

\subsection{Threshold memo template}
\begin{lstlisting}
\tau selection:
- Clinical prevalence: __
- FP/FN cost ratio: __
- Decision-curve net benefit at \tau: __
- Persistence rule effect (PPV \u0394, Sensitivity \u0394): __
Final \tau*: __   Rationale: __
\end{lstlisting}

\subsection{sql/cohort.sql (skeleton)}
\begin{lstlisting}[language=SQL]
WITH base AS (
  SELECT p.person_id, v.visit_date, v.hba1c, v.fpg, v.bmi, v.sbp, v.dbp, v.hrate,
         v.alt, v.ast, v.egfr, v.crp, v.ldl, v.hdl, v.tg,
         med.atyp_antipsych, med.steroids, med.statins,
         diag.icd10_code
  FROM visits v
  JOIN persons p   ON p.person_id = v.person_id
  LEFT JOIN meds med  ON med.person_id = v.person_id AND med.visit_id = v.visit_id
  LEFT JOIN diags diag ON diag.person_id = v.person_id AND diag.visit_id = v.visit_id
),
labels AS (
  SELECT person_id,
         MIN(CASE WHEN icd10_code LIKE 'E11%' THEN visit_date END) AS first_e11,
         MIN(CASE WHEN hba1c >= 6.5 THEN visit_date END) AS first_hba1c,
         MIN(CASE WHEN new_t2d_rx = 1 THEN visit_date END) AS first_rx
  FROM base
  GROUP BY person_id
)
SELECT *
FROM base b
LEFT JOIN labels l USING (person_id);
\end{lstlisting}

\section*{Acknowledgements}
We follow best practices for clinical ML evaluation, calibration, and reporting; any deployment must be preceded by institutional review and governance.

\begin{thebibliography}{99}
\bibitem[Brier(1950)]{brier1950verification} Brier, G. W. (1950). Verification of forecasts expressed in terms of probability. \emph{Monthly Weather Review}, 78(1), 1--3.
\bibitem[Chen and Guestrin(2016)]{chen2016xgboost} Chen, T., \& Guestrin, C. (2016). XGBoost: A scalable tree boosting system. In \emph{KDD}.
\bibitem[DeLong et al.(1988)]{delong1988comparing} DeLong, E. R., DeLong, D. M., \& Clarke-Pearson, D. L. (1988). Comparing the areas under two or more correlated ROC curves: A nonparametric approach. \emph{Biometrics}, 837--845.
\bibitem[Efron(1979)]{efron1979bootstrap} Efron, B. (1979). Bootstrap methods: Another look at the jackknife. \emph{Annals of Statistics}, 7(1), 1--26.
\bibitem[Kairouz et al.(2021)]{kairouz2021advances} Kairouz, P., et al. (2021). Advances and Open Problems in Federated Learning. \emph{Foundations and Trends in Machine Learning}, 14(1-2), 1--210.
\bibitem[Lin et al.(2017)]{lin2017focal} Lin, T.-Y., Goyal, P., Girshick, R., He, K., \& Doll\'ar, P. (2017). Focal Loss for Dense Object Detection. In \emph{ICCV}.
\bibitem[Loshchilov and Hutter(2019)]{loshchilov2019adamw} Loshchilov, I., \& Hutter, F. (2019). Decoupled weight decay regularization. In \emph{ICLR}.
\bibitem[Lundberg and Lee(2017)]{lundberg2017unified} Lundberg, S. M., \& Lee, S.-I. (2017). A Unified Approach to Interpreting Model Predictions. In \emph{NIPS}.
\bibitem[McNemar(1947)]{mcnemar1947note} McNemar, Q. (1947). Note on the sampling error of the difference between correlated proportions or percentages. \emph{Psychometrika}, 12(2), 153--157.
\bibitem[Platt(1999)]{platt1999probabilistic} Platt, J. (1999). Probabilistic outputs for SVMs and comparisons to regularized likelihood methods. \emph{Advances in Large Margin Classifiers}.
\bibitem[Spiegelhalter(1986)]{spiegelhalter1986probabilistic} Spiegelhalter, D. J. (1986). Probabilistic prediction in patient management and clinical trials. \emph{Statistics in Medicine}, 5(5), 421--433.
\bibitem[Sundararajan et al.(2017)]{sundararajan2017axiomatic} Sundararajan, M., Taly, A., \& Yan, Q. (2017). Axiomatic attribution for deep networks. In \emph{ICML}.
\bibitem[Vaswani et al.(2017)]{vaswani2017attention} Vaswani, A., et al. (2017). Attention Is All You Need. In \emph{NIPS}.
\bibitem[Vickers and Elkin(2006)]{vickers2006decision} Vickers, A. J., \& Elkin, E. B. (2006). Decision curve analysis: a novel method for evaluating prediction models. \emph{Medical Decision Making}, 26(6), 565--574.
\bibitem[Weinberger et al.(2009)]{weinberger2009feature} Weinberger, K., Dasgupta, A., Langford, J., Smola, A., \& Attenberg, J. (2009). Feature Hashing for Large Scale Multitask Learning. In \emph{ICML}.
\bibitem[Zadrozny and Elkan(2002)]{zadrozny2002transforming} Zadrozny, B., \& Elkan, C. (2002). Transforming classifier scores into accurate multiclass probability estimates. In \emph{KDD}.
\end{thebibliography}

\end{document}